\par\addvspace{12pt}\Needspace{9\baselineskip}

\section{切线、法线及其几何应用}

% ITEM: quadratic-example-070
\solutionitem{quadratic-example-070}{例4.70}
分别求椭圆\(x^2/9+y^2/4=1\)在\((3\sqrt2/2,\sqrt2)\)处、双曲线\(x^2/4-y^2=1\)在\((2\sqrt2,1)\)处及抛物线\(y^2=4x\)在\((1,2)\)处的切线和法线。
\solutionanswer
椭圆切线为\(2x+3y=6\sqrt2\)，法线可写为\((x,y)=(3\sqrt2/2,\sqrt2)+t(2,3)\)。双曲线切线为\(x/\sqrt2-y=1\)，法线\(y=-\sqrt2x+5\)。抛物线切线\(y=x+1\)，法线\(y=-x+3\)。每条法线都经过切点，方向取相应切线的法向量。


% ITEM: quadratic-example-071
\solutionitem{quadratic-example-071}{例4.71}
求椭圆\(x^2/9+y^2/4=1\)上切线平行于\(x+y=0\)的点。再将椭圆伸缩为单位圆，从圆上的相应点出发，核对所得切线。
\solutionanswer
切线法向量须平行\((1,1)\)，即\(x/9=y/4\)，与椭圆联立得\(P=\pm(9/\sqrt{13},4/\sqrt{13})\)。切线为\(x+y=\pm\sqrt{13}\)。圆化切点\(\pm(3/\sqrt{13},2/\sqrt{13})\)，切线\(3u+2v=\pm\sqrt{13}\)还原即得。


% ITEM: quadratic-example-072
\solutionitem{quadratic-example-072}{例4.72 *}
对曲线\(x^2+2xy+2y^2=1\)，用二次型与双线性表达求\((1,0)\)处的切线，再把切线参数代入曲线，核对交点的重数。
\solutionanswer
\(H=\begin{bmatrix}1&1\\1&2\end{bmatrix}\)，\(P=(1,0)^T\)，切线\(P^THX=1\)，即\(x+y=1\)。代入\(X=P+t(1,-1)\)，得\(q(X)-1=t^2\)，故交点参数\(0\)确为二重根。


% ITEM: quadratic-exercise-42-01
\solutionitem{quadratic-exercise-42-01}{习题4.42（1） *}
对\(y^2=4x\)在\((1,2)\)处的任意直线\(X=(1,2)+t(u,v)\)，写出限制多项式，并找出切线方向。
\par
\solutionanswer
限制式为\(4t(v-u)+v^2t^2\)。切线要求一次系数为零，故方向\(v=u\)，例如\((1,1)\)；此时\(t=0\)为二重根。


% ITEM: quadratic-exercise-42-02
\solutionitem{quadratic-exercise-42-02}{习题4.42（2） *}
对\(xy=1\)，分别研究直线\(y=x\)、\(y=2-x\)、\(x=1\)、\(x=0\)的限制多项式，比较次数、实根及重数。
\par
\solutionanswer
以横或纵坐标为参数，四个限制式依次为\(t^2-1,-(t-1)^2,t-1,-1\)。因此依次有两个单根、一个二重根、一个单根、无根。第三种只有一个有限实交点，却不是相切；它平行于渐近方向而使次数降低。


% ITEM: quadratic-exercise-42-03
\solutionitem{quadratic-exercise-42-03}{习题4.42（3） *}
设二次曲线非退化，\(P\)在曲线上。证明切线方向不可能同时满足\(q(V)=0\)，并用整条直线成为分量解释原因。
\par
\solutionanswer
\(P\)在曲线上使常数项为零，切线使一次项为零；若二次项\(q(V)\)也为零，则限制多项式恒为零，整条直线都属于曲线，成为一个直线分量，与非退化矛盾。


% ITEM: quadratic-exercise-42-04
\solutionitem{quadratic-exercise-42-04}{习题4.42（4） *}
若\(\Phi(P)\ne0\)，证明经过\(P\)、方向为\(V\)的直线在\(q(V)\ne0\)时相切的条件为
\(\bigl(2b(P,V)+L^TV\bigr)^2=4q(V)\Phi(P)\)。
\par
\solutionanswer
限制式为\(q(V)t^2+(2b(P,V)+L^TV)t+\Phi(P)\)。二次系数非零时，出现二重根当且仅当判别式为零，直接得到所求等式。


% ITEM: quadratic-exercise-42-05
\solutionitem{quadratic-exercise-42-05}{习题4.42（5） *}
对\(x^2+2xy+2y^2-2x=0\)，写出完整矩阵，求原点处切线，并用\(u=x+y-1,v=y\)将曲线化为圆，核对切线的变换。
\par
\solutionanswer
完整矩阵为\(\begin{bmatrix}1&1&-1\\1&2&0\\-1&0&0\end{bmatrix}\)。原点处一次变化为\(-2x\)，故切线\(x=0\)。给定代换得到\(u^2+(v+1)^2=2\)，原点映为\((-1,0)\)，该点切线为\(-u+v=1\)，还原正是\(x=0\)。


% ITEM: quadratic-example-073
\solutionitem{quadratic-example-073}{例4.73}
椭圆\(x^2/9+y^2/4=1\)与双曲线\(x^2-y^2/4=1\)有相同的焦点。求第一象限的交点及两条切线，并从坐标计算与焦点角平分关系两方面验证正交相交。
\solutionanswer
第一象限交点为\((3/\sqrt5,4/\sqrt5)\)。两切线分别为\(x+3y=3\sqrt5\)、\(3x-y=\sqrt5\)，斜率乘积为\(-1\)。从焦点看，椭圆切线平分两焦半径的外角，双曲线切线平分相应内角，两角平分线互相垂直，亦得正交。


% ITEM: quadratic-exercise-43-01
\solutionitem{quadratic-exercise-43-01}{习题4.43（1）}
求圆\((x-1)^2+(y+2)^2=5\)在\((2,0)\)处的切线和法线。说明圆的法线与圆心的关系。
\par
\solutionanswer
圆心为\(C=(1,-2)\)，半径向量\(\overrightarrow{CP}=(1,2)\)。切线垂直于半径，故
\[(1,2)\cdot(x-2,y)=0,\quad x+2y-2=0.\]
法线沿半径方向并经过\(P=(2,0)\)，可写为\((x,y)=(2,0)+s(1,2)\)，即\(y=2x-4\)。圆上任一点的半径都垂直于切线，因此圆的法线恒经过圆心；此处法线就是直线\(CP\)。


% ITEM: quadratic-exercise-43-02
\solutionitem{quadratic-exercise-43-02}{习题4.43（2）}
求抛物线\(y^2=8x\)上切线与\(x\)轴成\(45^\circ\)角的点，并求相应法线。
\par
\solutionanswer
将\(y^2=8x\)写作\(y^2=2px\)，有\(p=4\)。在\((x_0,y_0)\)处的切线为\(y_0y=4(x+x_0)\)，其斜率为\(4/y_0\)。与\(x\)轴成\(45^\circ\)的直线斜率为\(\pm1\)，故\(y_0=\pm4\)，而\(x_0=y_0^2/8=2\)。切点为\((2,4),(2,-4)\)。法线公式为\(y_0(x-x_0)+4(y-y_0)=0\)，分别得到
\[x+y-6=0,\qquad y-x+6=0.\]


% ITEM: quadratic-exercise-43-03
\solutionitem{quadratic-exercise-43-03}{习题4.43（3）}
设\(P\)在椭圆\(x^2/a^2+y^2/b^2=1\)上，\(a>b>0\)。证明法线经过原点当且仅当\(P\)为顶点。
\par
\solutionanswer
在 \(P=(u,v)\) 处，法线方向为 \((u/a^2,v/b^2)\)。它经过中心，当且仅当此方向与半径 \((u,v)\) 平行，即 \(uv(b^{-2}-a^{-2})=0\)。因 \(a>b\)，只有 \(u=0\) 或 \(v=0\)，在椭圆上恰给出四个顶点。


% ITEM: quadratic-exercise-43-04
\solutionitem{quadratic-exercise-43-04}{习题4.43（4）}
将圆\(u^2+v^2=1\)作变换\(x=2u+v,y=v\)。求像曲线在\((2,0)\)处的切线与法线，并比较它们与原来切线、法线的像。
\par
\solutionanswer
原圆上对应点为\((u,v)=(1,0)\)，其切线为\(u=1\)、法线为\(v=0\)。由逆变换\(u=(x-y)/2\)、\(v=y\)，像曲线方程为
\[
\frac{(x-y)^2}{4}+y^2=1.
\]
点\((2,0)\)处，切线的一个法向量为\((1,-1)\)，故像曲线切线为\(x-y=2\)，法线为\(x+y=2\)。切线也正是原切线\(u=1\)的像。原法线\(v=0\)的像是\(y=0\)，它与像曲线在\((2,0)\)处的法线不同；仿射变换保持切线作为割线极限的像，但一般不保持垂直关系，故法线不能简单随原法线取像。


% ITEM: quadratic-exercise-43-05
\solutionitem{quadratic-exercise-43-05}{习题4.43（5）}
椭圆\(x^2/25+y^2/9=1\)上一点\(P\)的横坐标为\(3\)。求两条焦半径的长度，并求两焦点到\(P\)处切线的距离之积。
\par
\solutionanswer
由 \(a=5,b=3,c=4\)，焦半径公式直接给出到左、右焦点的距离 \(5\pm(4/5)\cdot3\)，分别为 \(37/5,13/5\)。求切线距离之积时也可以将两焦点配成一对：由椭圆方程 \(y_P=\pm12/5\)，相应切线为 \(9x\pm20y=75\)。对称的两个焦点代入后，只改变 \(9x\) 的符号，故乘积为
\[
 \frac{(75-36)(75+36)}{9^2+20^2}=9.
\]
这个值恰为 \(b^2\)。下一题说明，换到别的切点也有同样的乘积。


% ITEM: quadratic-exercise-43-06
\solutionitem{quadratic-exercise-43-06}{习题4.43（6）}
证明椭圆两焦点到任意切线的距离之积为\(b^2\)。可以利用切线的单位法向量，也可以利用焦点的反射点。
\par
\solutionanswer
取切线的单位法向量 \(n=(n_x,n_y)\)，使方程为 \(n\cdot X=h>0\)。将椭圆写成 \((a\cos t,b\sin t)\)，可知 \(h=\max_t(an_x\cos t+bn_y\sin t)=\sqrt{a^2n_x^2+b^2n_y^2}\)。两个焦点都在切线同侧，到切线的距离为 \(h-cn_x,\ h+cn_x\)。因此
\[
 d_1d_2=h^2-c^2n_x^2=b^2(n_x^2+n_y^2)=b^2.
\]

也可直接从焦点三角形看这个乘积。设两焦半径为 \(r_1,r_2\)，夹角为 \(\varphi\)。切线平分两焦半径的外角，所以两焦点到切线的距离分别为 \(r_1\cos(\varphi/2),r_2\cos(\varphi/2)\)。对边是固定的焦距 \(2c\)，由余弦定理，
\[
 4d_1d_2=2r_1r_2(1+\cos\varphi)=(r_1+r_2)^2-(2c)^2=4a^2-4c^2.
\]
在长轴顶点，两距离直接为 \(a-c,a+c\)，乘积相同。


% ITEM: quadratic-exercise-43-07
\solutionitem{quadratic-exercise-43-07}{习题4.43（7）}
求抛物线\(y^2=8x\)的焦点在各条切线上的垂足轨迹。若垂足为\((0,2)\)，求切点。
\par
\solutionanswer
用切点参数\((2t^2,4t)\)，切线为\(ty=x+2t^2\)。焦点\((2,0)\)在此线上的垂足为\((0,2t)\)，故轨迹为整个\(y\)轴。给定垂足\((0,2)\)对应\(t=1\)，切点为\((2,4)\)。


% ITEM: quadratic-exercise-43-08
\solutionitem{quadratic-exercise-43-08}{习题4.43（8）}
固定\(F_1,F_2\)，取\(2a>|F_1F_2|\)。在以\(F_2\)为圆心、\(2a\)为半径的圆上取\(S\)，将\(S\)折到\(F_1\)。证明折痕与线段\(F_2S\)的交点属于以\(F_1,F_2\)为焦点的椭圆，且折痕是该点处的切线。
\par
\solutionanswer
折痕是\(F_1S\)的垂直平分线。因\(F_2F_1<2a=F_2S\)，\(F_2\)在靠近\(F_1\)的一侧，\(S\)在另一侧，故线段\(F_2S\)与折痕恰交于一点\(P\)。由\(PF_1=PS\)得\(PF_1+PF_2=PS+PF_2=2a\)，在椭圆上。折痕上任一点\(X\)满足\(XF_1+XF_2=XS+XF_2\ge SF_2=2a\)，等号只在\(P\)取得。因此折痕是椭圆在\(P\)处的支持线，即切线。


% ITEM: quadratic-exercise-43-09
\solutionitem{quadratic-exercise-43-09}{习题4.43（9）}
证明双曲线\(x^2/a^2-y^2/b^2=1\)的一个焦点在各条切线上的垂足均在圆\(x^2+y^2=a^2\)上。这个圆是否被全部覆盖？将未覆盖的位置与渐近线联系起来。
\par
\solutionanswer
取两焦点 \(F=(c,0),G=(-c,0)\)，将 \(F\) 关于切点 \(P\) 处的切线反射到 \(U\)。由双曲线的反射性质，\(U\) 在射线 \(PG\) 上，故 \(GU=|PG-PF|=2a\)。垂足 \(T\) 是 \(FU\) 的中点，而原点 \(O\) 是 \(FG\) 的中点，于是 \(\overrightarrow{OT}=\tfrac12\overrightarrow{GU}\)，立即得到 \(OT=a\)。

反过来，取圆上点 \(T\)，作 \(F\) 关于 \(T\) 的对称点 \(U\)，则 \(GU=2a\)。候选切线是 \(FU\) 的中垂线 \(l\)。因 \(GF=2c>GU\)，\(G,U\) 在 \(l\) 同侧。若 \(l\) 与直线 \(GU\) 相交，交点 \(P\) 在线段 \(GU\) 外，且 \(PF=PU\)，故 \(|PG-PF|=GU=2a\)。折痕 \(l\) 将射线 \(PF\) 与 \(PG\) 互换，正是该点处的切线（在顶点处为垂直于实轴的直线）。

只剩 \(l\parallel GU\) 的情形。此时 \(l\) 就是 \(OT\)，故 \(FT\perp OT\)。写 \(T=(u,v)\)，即 \(cu=u^2+v^2=a^2\)。圆上恰漏掉 \((a^2/c,\pm ab/c)\)，相应的 \(OT\) 正是两条渐近线 \(y=\pm(b/a)x\)。


% ITEM: quadratic-exercise-43-10
\solutionitem{quadratic-exercise-43-10}{习题4.43（10）}
设\(a>b>0,\lambda>0\)。证明椭圆\(x^2/a^2+y^2/b^2=1\)与双曲线\(x^2/(a^2-\lambda)-y^2/(\lambda-b^2)=1\)在\(b^2<\lambda<a^2\)时共焦，并求证其交点处切线互相垂直。
\par
\solutionanswer
双曲线的焦距平方为\((a^2-\lambda)+(\lambda-b^2)=a^2-b^2\)，故共焦。两曲线方程相减后除以\(\lambda\)，得到\(x^2/[a^2(a^2-\lambda)]-y^2/[b^2(\lambda-b^2)]=0\)。左边正是两切线法向量\((x/a^2,y/b^2)\)、\((x/(a^2-\lambda),-y/(\lambda-b^2))\)的数量积，因此两切线垂直。


% ITEM: quadratic-example-074
\solutionitem{quadratic-example-074}{例4.74}
（2025年天津卷第18题）椭圆\(x^2/a^2+y^2/b^2=1\)的左焦点为\(F\)，右顶点为\(A\)，离心率为\(1/2\)。点\(P\)在\(x=a\)上，\(PF\)斜率为\(1/3\)，且\(S_{\triangle PFA}=3/2\)。求椭圆方程；过\(P\)的另一条切线切于\(B\)，证明\(FP\)平分\(\angle AFB\)。
\solutionanswer
由 \(c=a/2\) 及 \(PF\) 的斜率得 \(P=(a,a/2)\)，故 \(3a^2/8=S_{PFA}=3/2\)。于是 \(a=2,c=1\)，椭圆为 \(x^2/4+y^2/3=1\)。

此时 \(P=(2,1)\)，切点弦为 \(x/2+y/3=1\)，与椭圆除 \(A\) 外还交于 \(B=(1,3/2)\)。射线 \(FA,FB\) 的单位向量之和为 \((1,0)+(4/5,3/5)=\tfrac35(3,1)\)，与 \(\overrightarrow{FP}\) 同向，故 \(FP\) 平分 \(\angle AFB\)。

另证不必求 \(B\)。把右焦点 \(G\) 分别关于 \(PA,PB\) 反射到 \(U,V\)。由反射性质，\(U,V\) 分别在射线 \(FA,FB\) 上，且 \(FU=FV=2a\)；又 \(PU=PV=PG\)，故 \(FP\) 是 \(UV\) 的中垂线。射线 \(FP\) 与 \(FA\) 成锐角，因而平分等腰三角形 \(FUV\) 的内角。


% ITEM: quadratic-example-075
\solutionitem{quadratic-example-075}{例4.75}
设\(D\)在直线\(y=-1/2\)上，过\(D\)作抛物线\(C:y=x^2/2\)的两条切线，切点为\(A,B\)。

（1）证明\(AB\)经过定点；

（2）若以\(E=(0,5/2)\)为圆心的圆在\(AB\)的中点处与\(AB\)相切，求该圆的方程。

（2019年全国卷III文科高考题改编）
\solutionanswer
（1）设\(D=(d,-1/2)\)，切点横坐标\(u,v\)满足\(t^2-2dt-1=0\)。故\(u+v=2d,uv=-1\)，切点弦\(y=dx+1/2\)，经过\((0,1/2)\)。（2）中点\(M=(d,d^2+1/2)\)。半径\(EM\)垂直弦要求\((d,d^2-2)\cdot(1,d)=d(d^2-1)=0\)，故\(d=0,\pm1\)。对应全部圆为\(x^2+(y-5/2)^2=4\)或\(x^2+(y-5/2)^2=2\)。


% ITEM: quadratic-exercise-44-01
\solutionitem{quadratic-exercise-44-01}{习题4.44（1）}
从椭圆外一点\(P=(u,v)\)作两条切线，切点为\(A,B\)，右焦点为\(F=(c,0)\)。利用切点弦与焦半径公式证明\(\displaystyle (P-F)\cdot\frac{A-F}{|AF|}=a-\frac{cu}{a}\)，再将\(A\)换成\(B\)，说明\(\angle AFP=\angle BFP\)。
\par
\solutionanswer
切点弦关系给出\(ux_A/a^2+vy_A/b^2=1\)，焦半径为\(a-cx_A/a\)。因此
\[
(P-F)\cdot(A-F)=(a-cx_A/a)(a-cu/a).
\]
相除得到题中等式。换成\(B\)右边相同，再除以\(|PF|\)，两个角的余弦相同，故两角相等。


% ITEM: quadratic-exercise-44-02
\solutionitem{quadratic-exercise-44-02}{习题4.44（2）}
从\(P=(0,3)\)向椭圆\(x^2/9+y^2/4=1\)作切线，求切点弦方程及弦长。
\par
\solutionanswer
切点弦为\(3y/4=1\)，即\(y=4/3\)。代入椭圆得端点\((\pm\sqrt5,4/3)\)，弦长\(2\sqrt5\)。


% ITEM: quadratic-exercise-44-03
\solutionitem{quadratic-exercise-44-03}{习题4.44（3）}
求过\((-1,2)\)且与抛物线\(y^2=4x\)相切的两条直线，并求切点弦方程。
\par
\solutionanswer
抛物线非竖直切线为\(y=kx+1/k\)。代入给定点得\(k^2+2k-1=0\)，故两斜率为\(-1\pm\sqrt2\)，相应两切线为\(y=kx+1/k\)。切点弦由观察点\((-1,2)\)得到\(2y=2(x-1)\)，即\(y=x-1\)。


% ITEM: quadratic-exercise-44-04
\solutionitem{quadratic-exercise-44-04}{习题4.44（4）}
求\(xy=1\)在\((t,1/t)\)处的切线，完整分类讨论过\((u,v)\)的实切线条数，并用渐近线方向解释次数降低的情形。
\par
\solutionanswer
切线为\(x/t+ty=2\)，\(t\ne0\)。过\((u,v)\)要求\(vt^2-2t+u=0\)，且须去掉\(t=0\)。因此：\(uv\ne0\)时，\(uv<1,uv=1,uv>1\)分别有\(2,1,0\)条；恰一个坐标为零时有\(1\)条；原点有\(0\)条。坐标为零时，一个形式根为零或移到无穷，对应渐近方向而不是有限切点，故条数降低。


% ITEM: quadratic-exercise-44-05
\solutionitem{quadratic-exercise-44-05}{习题4.44（5）}
设\(P\)在椭圆外，\(PA\)为切线，另一条直线交椭圆于\(B,C\)，相应单位方向为\(d_1,d_2\)。记\(q(u,v)=u^2/a^2+v^2/b^2\)，证明\(|PA|^2q(d_1)=\overline{PB}\,\overline{PC}\,q(d_2)\)，并说明何时能得到通常的切割线长度关系。
\par
\solutionanswer
两种限制式的常数均为\(q(P)-1\)。切线的重根给出\(|PA|^2q(d_1)=q(P)-1\)，割线的根积给出\(\overline{PB}\,\overline{PC}q(d_2)=q(P)-1\)，故等式成立。普通切割线长度关系需要两方向的\(q\)值相等；圆中所有单位方向的\(q\)值相同，因此总成立。


% ITEM: quadratic-exercise-44-06
\solutionitem{quadratic-exercise-44-06}{习题4.44（6）}
抛物线\(C:x^2=2py\)的焦点为\(F\)，圆\(M:x^2+(y+4)^2=1\)。已知\(F\)到圆\(M\)上各点的距离最小值为\(4\)。求\(p\)；若\(P\)在圆\(M\)上，过\(P\)作\(C\)的两条切线，切点为\(A,B\)，求\(\triangle PAB\)面积的最大值。（2021年全国乙卷理科高考题改编）
\par
\solutionanswer
焦点为 \((0,p/2)\)，到圆心的距离为 \(4+p/2\)，减去半径得 \(3+p/2=4\)，故 \(p=2\)。在 \(x^2=4y\) 上取切点 \(T_i=(2t_i,t_i^2)\)，其切线为 \(y=t_ix-t_i^2\)。两切线相交于 \(P=(u,v)\)，故 \(t_1,t_2\) 是 \(t^2-ut+v=0\) 的两根，即 \(P=(t_1+t_2,t_1t_2)\)。从 \(P\) 指向两切点的向量为
\[
 T_1-P=(t_1-t_2)(1,t_1),\qquad T_2-P=(t_2-t_1)(1,t_2).
\]
取行列式便得 \(2S=|t_1-t_2|^3=(u^2-4v)^{3/2}\)。令 \(z=v+4\)，则 \(u^2+z^2=1\)，故 \(u^2-4v=21-(z+2)^2\)。因 \(-1\le z\le1\)，此式在 \(z=-1\) 时最大，为 \(20\)。于是最大面积为 \(20\sqrt5\)，此时 \(P=(0,-5)\)。


% ITEM: quadratic-exercise-44-07
\solutionitem{quadratic-exercise-44-07}{习题4.44（7）}
椭圆\(E:x^2/a^2+y^2/b^2=1\)的左、右焦点为\(F_1,F_2\)。过\(F_1\)的竖直直线与上半椭圆交于\(P\)，过\(F_2\)且垂直于\(PF_2\)的直线与右准线交于\(Q\)。若\(Q=(4,4)\)，求\(E\)的方程，并证明\(PQ\)是其切线。（2012年安徽理科高考题改编）
\par
\solutionanswer
由右准线得 \(a^2=4c\)，所以 \(b^2=c(4-c)\)。又 \(P=(-c,b^2/a)\)，由 \(F_2P\perp F_2Q\) 得
\[
 0=(-2c,b^2/a)\cdot(4-c,4)=-2b^2+4b^2/a.
\]
故 \(a=2,\ c=1,\ b^2=3\)，椭圆为 \(x^2/4+y^2/3=1\)。在 \(P=(-1,3/2)\) 处的切线为 \(-x/4+y/2=1\)，而 \(Q=(4,4)\) 满足此式，故 \(PQ\) 即切线。

也可直接从作图看出相切。写 \(P=(u,v),\ Q=(a^2/c,w)\)，则垂直条件为 \((u-c)b^2/c+vw=0\)，即 \(u/c+vw/b^2=1\)。这正是 \(Q\) 满足 \(P\) 处切线方程 \(ux/a^2+vy/b^2=1\) 的条件。


% ITEM: quadratic-example-076
\solutionitem{quadratic-example-076}{例4.76}
过椭圆\(x^2/9+y^2/4=1\)的两个不同点分别作切线，两切线互相垂直且交点在\(y=x\)上。求所有交点及相应的切点弦。
\solutionanswer
设交点为 \(T=(u,u)\)。由准圆方程 \(x^2+y^2=9+4\)，得 \(u=\pm\sqrt{13/2}\)。相应的切点弦是 \(T\) 的极线 \(ux/9+uy/4=1\)，即 \(4x+9y=36/u\)。

另看准圆的几何来历。从右焦点 \(F\) 向两切线作垂线，垂足为 \(H,K\)，则 \(OH=OK=3\)，且 \(FHTK\) 是矩形。两条对角线等长且互相平分，对它们各用一次中线公式，即得 \(OF^2+OT^2=OH^2+OK^2\)，所以 \(OT^2=18-5=13\)。


% ITEM: quadratic-example-077
\solutionitem{quadratic-example-077}{例4.77}
双曲线\(x^2/9-y^2/4=1\)在\(P\)处的切线与两渐近线交于\(A,B\)。证明\(P\)是\(AB\)的中点，求\(\triangle OAB\)的面积，并求\(|AB|\)的最小值。
\solutionanswer
既然要求切线在两条渐近线上的截点，就取 \(u=x/3+y/2,\ v=x/3-y/2\)。双曲线成为 \(uv=1\)，两条渐近线成为坐标轴。设切点的新坐标为 \((s,t)\)，\(st=1\)，切线为 \(tu+sv=2\)：将 \(v=(2-tu)/s\) 代入 \(uv=1\)，所得方程正是 \(t(u-s)^2=0\)。两截点为 \((2s,0),(0,2t)\)，其中点就是 \((s,t)\)。线性变换保持中点，故原图中 \(P\) 平分 \(AB\)。

面积和长度回到原坐标计算。由 \(x=3(u+v)/2,\ y=u-v\)，得 \(A=(3s,2s), B=(3t,-2t)\)，故
\[
 S_{OAB}=6|st|=6,\qquad |AB|^2=9(s-t)^2+4(s+t)^2=16+13(s-t)^2.
\]
所以截得的线段 \(AB\) 的最小长度为 \(4\)，在 \(s=t=\pm1\)，即 \(P\) 为任一实轴顶点时取得。


% ITEM: quadratic-example-078
\solutionitem{quadratic-example-078}{例4.78}
求抛物线\(y^2=4x\)与直线\(x=4\)围成的面积。若将直线改为\(x-y=2\)，仍用平行于弦的切线确定内接三角形，再求弓形面积。
\solutionanswer
弓形面积为相应内接三角形面积的 \(4/3\)。对竖直弦 \(x=4\)，端点为 \((4,\pm4)\)，平行切线在顶点，三角形的底为 \(8\)、高为 \(4\)，故弓形面积为 \(64/3\)。

对斜弦，将 \(x-y=y^2/4-y=(y-2)^2/4-1\) 配方。取 \(u=y-2,\ v=x-y+1\)，面积倍率为 \(1\)。此时抛物线为 \(v=u^2/4\)，弦为 \(v=3\)，平行切线为 \(v=0\)（原切点为 \((1,2)\)）。弦长在新坐标中是 \(4\sqrt3\)，高为 \(3\)，故原弓形面积为 \(\tfrac43\cdot\tfrac12\cdot4\sqrt3\cdot3=8\sqrt3\)。


% ITEM: quadratic-exercise-45-01
\solutionitem{quadratic-exercise-45-01}{习题4.45（1）}
求椭圆\(x^2/16+y^2/9=1\)的准圆。若两条垂直切线的交点在\(x\)轴上，求切点弦长。
\par
\solutionanswer
准圆为\(x^2+y^2=25\)，与\(x\)轴交于\((\pm5,0)\)。切点弦为\(x=\pm16/5\)，端点纵坐标\(\pm9/5\)，故弦长\(18/5\)。


% ITEM: quadratic-exercise-45-02
\solutionitem{quadratic-exercise-45-02}{习题4.45（2）}
设双曲线为\(x^2/a^2-y^2/b^2=1\)。仿照单位法向量的方法，研究两条互相垂直的实切线的交点轨迹，并比较\(a>b\)、\(a=b\)、\(a<b\)三种情形。注意哪些位置不能产生两条实切线。
\par
\solutionanswer
设交点为 \(P=(u,v)\)，过它的直线取单位法向量 \((m,n)\)，写成 \(mx+ny=h=um+vn\)。与标准切线 \(x_0x/a^2-y_0y/b^2=1\) 比较，有限相切要求 \(h^2=a^2m^2-b^2n^2,\ h\ne0\)，切点为 \((a^2m/h,-b^2n/h)\)。若等式成立而 \(h=0\)，直线是渐近线。

将上述等式写为 \(D(m,n)=(a^2-u^2)m^2-2uvmn-(b^2+v^2)n^2=0\)。两个法向量垂直时，可取为 \((m,n),(-n,m)\)；两式相加即得 \(u^2+v^2=a^2-b^2\)。反过来，在这个圆上令 \((m,n)=(\cos\theta,\sin\theta)\)，条件简化成
\[
 (b^2+v^2)\cos 2\theta-uv\sin 2\theta=0 .
\]
它恰给出相差 \(\pi/2\) 的两个法向。当 \(a>b\) 时，只需去掉圆与两条渐近线的四个交点，因为那里有一个方向不给出有限切点。当 \(a=b\) 时仅剩原点，两方向都是渐近方向；\(a<b\) 时无实圆。这两种情形的轨迹都为空。


% ITEM: quadratic-exercise-45-03
\solutionitem{quadratic-exercise-45-03}{习题4.45（3）}
求双曲线\(xy=3\)在\((t,3/t)\)处的切线与两渐近线所成三角形面积，并求其周长的最小值。
\par
\solutionanswer
切线为\(3x/t+ty=6\)，两轴截距\(2t,6/t\)，三角形面积恒为\(6\)。设两直角边长\(r=2|t|,s=6/|t|\)，则\(rs=12\)。由\(r+s\ge4\sqrt3\)、\(\sqrt{r^2+s^2}\ge2\sqrt6\)，周长最小\(4\sqrt3+2\sqrt6\)，在\(|t|=\sqrt3\)取到。


% ITEM: quadratic-exercise-45-04
\solutionitem{quadratic-exercise-45-04}{习题4.45（4）}
证明任意外切于椭圆\(x^2/a^2+y^2/b^2=1\)的平行四边形面积不小于\(4ab\)，并说明等号与共轭直径的关系。提示：先变成单位圆的外切平行四边形。
\par
\solutionanswer
圆化后是单位圆的外切平行四边形，两对平行边间距均为\(2\)。若相邻边夹角为\(\theta\)，面积为\(4/|\sin\theta|\ge4\)，等号为正方形。还原后面积不小于\(4ab\)，等号对应四边在一对共轭直径端点处相切。


% ITEM: quadratic-exercise-45-05
\solutionitem{quadratic-exercise-45-05}{习题4.45（5）}
设\(y=x^2\)上弦端点的横坐标之差为\(d>0\)。求弓形面积，说明它是否取决于弦的斜率。再求\(y=\alpha x^2+\beta x+\gamma\)中相应的面积公式，其中\(\alpha\ne0\)。
\par
\solutionanswer
端点横坐标差为\(d\)时，平行于弦的切点位于两横坐标的平均处，弦到曲线的竖直差为\(d^2/4\)。内接三角形面积\(d^3/8\)，弓形面积为其\(4/3\)，即\(d^3/6\)，与斜率无关。一般二次函数通过平移、错切及竖直伸缩还原，面积为\(|\alpha|d^3/6\)。


% ITEM: quadratic-exercise-45-06
\solutionitem{quadratic-exercise-45-06}{习题4.45（6）}
抛物线\(y^2=4x\)在\(A,B\)处的两条切线交于\(P\)。证明弦\(AB\)与抛物线围成的弓形面积为\(\triangle PAB\)面积的\(2/3\)。
\par
\solutionanswer
在抛物线上取参数点 \((t^2,2t)\)。将弦的端点参数写为 \(s-d,s+d\)（\(d>0\)），则弦中点为 \(M=(s^2+d^2,2s)\)。由端点切线 \(x=ty-t^2\)，两切线交于 \(P=(s^2-d^2,2s)\)。平行于弦的切线切于 \(T=(s^2,2s)\)：弦和平行切线的方向均为 \((s,1)\)。

由此 \(T\) 恰为 \(PM\) 的中点。以 \(AB\) 为同一底边，\(T\) 的高就是 \(P\) 的高的一半。抛物线弓形的面积为 \(S_{TAB}\) 的 \(4/3\)，因而正是 \(S_{PAB}\) 的 \(2/3\)。
\par
